\begin{question}
\noindent The diagram shows a regular nonagon (a 9-sided polygon).

\begin{center}
\begin{tikzpicture}[scale=1.3]
    % Regular nonagon with 9 sides, interior angle marked at one vertex
    \def\n{9}
    \def\R{2}
    \foreach \i in {0,...,8} {
        \pgfmathsetmacro{\angA}{360/\n * \i + 90}
        \pgfmathsetmacro{\angB}{360/\n * (\i+1) + 90}
        \coordinate (P\i) at (\angA:\R);
    }
    \draw[thick] (P0) -- (P1) -- (P2) -- (P3) -- (P4) -- (P5) -- (P6) -- (P7) -- (P8) -- cycle;

    \pic [draw, "\footnotesize $x$", angle eccentricity=1.6, angle radius=0.5cm] {angle = P1--P0--P8};

    \node at (0,-2.6) {\footnotesize Regular nonagon (9 equal sides, 9 equal angles)};
\end{tikzpicture}
\end{center}

\noindent The diagram shows a regular nonagon (9 sides).

\noindent Work out the size of the interior angle marked $x$.
\vspace{4.5cm}
\answerequnit{x}{$^\circ$}{3}
\end{question}
